\section{Theoretical Analysis}
\label{sec:theory}

% 证据：实验 6a（theory/model.py、sweep.py、sweep.json、meanfield.py）；T11 报告（bias_multi.py，检验 B、D；t11_tables.md）
% 不能写："理论预言了 K0.5"；区间比例 1/q 在 FJSP 上未定量验证（只在实验节写定性结论）
% 2026-10-22 命题 1～3、推论、K* = 12/26/55、表格数字均已独立核对；本版只改写表达，事实未改。
% 待补（T12 结果出来后再写）：冗余解码是否降低多余偏离的代价

How many deviation slots should \RDD use?
Too few, and the search cannot reach rosters that differ from the rule in many shifts.
Too many, and random genotypes contain many unnecessary deviations, so the search starts far from the rule and wastes evaluations removing them.
This section makes the trade-off precise with a simple model.
The model shows that the best number of slots depends on two quantities:
how many decisions separate a good solution from the rule, and how many evaluations the search can spend.

\subsection{Model}
\label{sec:theory:model}

A solution is a string $\mathbf{z} \in \{0, \dots, \nval-1\}^{\npos}$ of $\npos$ decisions with $\nval$ options each.
Option 0 is the rule's choice, so the rule's solution is the all-zero string.
A target solution $\mathbf{z}^\ast$ differs from the rule in $\dist$ positions.
We call $\dist$ the \emph{rule distance}.
The objective is the number of positions in which a solution differs from the target, $f(\mathbf{z}) = |\{j : z_j \ne z^\ast_j\}|$.
In this model, changing one decision does not change the meaning of other decisions.
Section~\ref{sec:motivation} found that this kind of propagation did not explain performance differences between decoders.

We compare two encodings with genes in $[0,1)$:
\begin{itemize}
  \item \textbf{Per-shift (PS)}: one gene per decision, read as $z_j = \lfloor g_j \nval \rfloor$.
  This is a simplified, unbiased version of \dec{0}--\dec{5}; the real decoders lean toward the rule's choice more than PS does.
  \item \textbf{Sparse deviation (SP$_\K$)}: the model version of \RDD, with $\K$ slots as in Section~\ref{sec:method:genotype}.
  A slot is on with probability $\qon \approx 0.4$ in a random genotype.
  An active slot sets one position, chosen by $p_k$, to a non-zero option chosen by $a_k$.
  If two slots hit the same position, the later one counts.
\end{itemize}

\subsection{Where the Search Starts and What It Can Reach}
\label{sec:theory:props}

We first look at three properties of the encoding itself, before any search.
The first is how far a random genotype is from the target.

\begin{proposition}[Starting distance]
\label{prop:init}
For a uniformly random genotype, let $\pi = 1 - (1 - \qon/\npos)^\K$ be the probability that a given position is changed by some slot. Then
\begin{equation}
  \mathbb{E}[f_{\mathrm{SP}}] = (\npos - \dist)\,\pi + \dist \left(1 - \frac{\pi}{\nval-1}\right).
\end{equation}
For PS, $\mathbb{E}[f_{\mathrm{PS}}] = \npos (1 - 1/\nval)$.
\end{proposition}

\begin{IEEEproof}
Each slot hits a given position with probability $\qon/\npos$, independently of the other slots, so the position is changed with probability $\pi$.
The last write at a changed position is uniform over the $\nval-1$ non-zero options.
A position where the target agrees with the rule is wrong exactly when it is changed.
A position where the target differs from the rule is right exactly when it is changed to the right option, which has probability $\pi/(\nval-1)$.
\end{IEEEproof}

When $\K$ and $\dist$ are small compared with $\npos$, this is about $\dist + \qon\K$, which depends on $\K$ and $\qon$ only through $\qon\K$, the expected number of active slots:
a random genotype is about as far from the target as the rule is, plus roughly one wrong change per active slot.
Under PS, a random genotype is $\npos(1-1/\nval)$ away, whatever the rule is; for $\npos = 100$ and $\nval = 4$ this is 75.

The second property is a hard limit.

\begin{proposition}[Reach]
\label{prop:floor}
A genotype of SP$_\K$ changes at most $\K$ positions, so no genotype gets closer to the target than $\max(0, \dist - \K)$.
\end{proposition}

With fewer slots than the rule distance, the target cannot be reached, however long the search runs.

The third property counts how many genotypes produce the target.
The more genotypes that do, the easier the target is to find.

\begin{proposition}[How often the target is represented]
\label{prop:red}
The probability that a uniformly random genotype of SP$_\K$ produces exactly $\mathbf{z}^\ast$ is
\begin{equation}
  r(\K) = \frac{\dist!}{(\nval-1)^{\dist}} \sum_{J=\dist}^{\K} \binom{\K}{J} \qon^J (1-\qon)^{\K-J} \npos^{-J}\, \mathcal{S}(J, \dist),
\end{equation}
where $\mathcal{S}(J,\dist)$ is the Stirling number of the second kind. In particular, $r(\K) = 0$ for $\K < \dist$.
\end{proposition}

\begin{IEEEproof}
Let $J$ be the number of active slots; $J$ is binomial with parameters $\K$ and $\qon$.
The genotype produces $\mathbf{z}^\ast$ exactly when all $J$ active slots fall on the $\dist$ target positions,
every target position is hit, and the last write at each one is the right option.
Of the $\npos^J$ equally likely placements, $\dist!\,\mathcal{S}(J,\dist)$ cover all $\dist$ positions.
Each of the $\dist$ last writes is right with probability $1/(\nval-1)$.
\end{IEEEproof}

\begin{corollary}
\label{cor:kstar}
Keeping only the leading term $J = \dist$, $r(\K+1)/r(\K) = (\K+1)(1-\qon)/(\K+1-\dist)$, which is at least one exactly when $\K + 1 \le \dist/\qon$.
The target is therefore represented most often when $\K$ is close to $\K^\ast \approx \dist/\qon$.
\end{corollary}

The exact maximizers agree.
For $\npos = 100$, $\nval = 4$ and $\qon = 0.4$, they are $\K^\ast = 12$, 26 and 55 for $\dist = 5$, 10 and 20, against $\dist/\qon = 12.5$, 25 and 50.
Under PS, the target is produced with probability $\nval^{-\npos}$, independent of the rule.

In short: with $\K < \dist$ the target is out of reach;
near $\K = \dist/\qon$ it is produced most often;
with more slots, random genotypes make too many wrong changes, and both the starting distance and the chance of producing the target get worse.

\subsection{The Best Number of Slots for a Given Budget}
\label{sec:theory:budget}

The three properties describe the encoding, not the search.
To find the best $\K$ for a given number of evaluations, we run a GA with the same operators as our NSGA-II (Section~\ref{sec:exp-setup}):
simulated binary crossover (SBX, $\eta = 20$, probability 0.9), polynomial mutation ($\eta = 20$, rate $1/n$ for a genotype of length $n$), population 30 and truncation selection.
We use $\npos = 100$ and $\nval = 4$, nine rule distances between 0 and 80, eleven values of $\K$ between 5 and 100, and 20 runs per setting.

Table~\ref{tab:theory-budget} shows that the best $\K$ grows with the budget:
\begin{itemize}
  \item With 465 evaluations and $\dist \le 40$, the best $\K$ is at most $\dist$.
  \item With 2000 evaluations, it is larger than $\dist$, and close to $\dist/\qon$ when $\dist$ is small.
  \item With 10\,000 evaluations, it is larger than $\dist/\qon$: 20, 50 and 75 for $\dist = 5$, 10 and 20.
\end{itemize}
No single value of $\K$ is good for all budgets.
For example, $\K = 25$ is the best choice in only one of the twelve settings of the table;
with $\dist = 40$ and 10\,000 evaluations it ends 29.7 positions from the target, while $\K = 100$ ends 12.1 away.
PS is far behind whenever $\dist \le 60$: it is still more than 40 positions away after 10\,000 evaluations.

\begin{table}[t]
  \centering
  \caption{Best number of slots $\K$ in the model ($\npos = 100$, $\nval = 4$, $\qon \approx 0.4$). Entries: best $\K$ (final distance to the target), mean of 20 runs.}
  \label{tab:theory-budget}
  \setlength{\tabcolsep}{4pt}
  \begin{tabular}{@{}rrccc@{}}
    \toprule
    $\dist$ & $\dist/\qon$ & 465 evals & 2000 evals & 10\,000 evals \\
    \midrule
    5  & 12.5 & 5 (4.0)       & 15 (3.5)  & 20 (2.4) \\
    10 & 25   & 5, 10 (8.6)   & 25 (6.8)  & 50 (3.8) \\
    20 & 50   & 10 (17.2)     & 30 (13.6) & 75 (6.8) \\
    40 & 100  & 20, 30 (35.1) & 60 (27.1) & 100$^\dagger$ (12.1) \\
    \bottomrule
    \multicolumn{5}{@{}l}{\footnotesize $^\dagger$Largest value tested.}
  \end{tabular}
\end{table}

\subsection{Why More Evaluations Call for More Slots}
\label{sec:theory:meanfield}

A simple model of the search dynamics explains this trend.
At any time, each slot is in one of three states:
on and wrong (it changes a position that should not change, or sets a wrong option), on and right, or off.
Let $w$, $c$ and $v = \K - w - c$ be the numbers of slots in these states, so that $f = w + (\dist - c)$.
Mutation changes about one gene per offspring, so each slot is touched with probability about $1/\K$.
Selection keeps only changes that lower $f$.
There are three such changes: switching off a wrong slot ($w \to w-1$),
switching on an inactive slot at a missing position with the right option ($c \to c+1$),
and moving a wrong slot to a missing position with the right option ($w \to w-1$, $c \to c+1$).
They occur at the rates
\begin{equation}
  \alpha\, \frac{w}{2\K}, \qquad
  \alpha\, \frac{v}{2\K} \cdot \frac{\dist - c}{\npos(\nval-1)}, \qquad
  \beta\, \frac{w}{2\K} \cdot \frac{\dist - c}{\npos(\nval-1)}.
\end{equation}
The starting values of $w$ and $c$ come from Proposition~\ref{prop:init}, corrected for selecting the best of the initial population.
The constants $\alpha$ and $\beta$ depend only on the operators.
We fit them on the runs with 465 and 2000 evaluations ($\alpha = 0.05$, $\beta = 1.5$; mean squared error 2.4).
On the runs with 10\,000 evaluations, which were not used for fitting, predicted and observed final distances have a correlation of 0.994,
but a larger mean squared error of 22.6, and the model's best $\K$ is 1.3 to 3 times too large.
The model therefore captures the trend but not the exact values; we use it only to explain the trend.

The explanation is as follows.
Every wrong slot costs one unit of distance and has to be switched off.
A random genotype starts with about $\qon\K$ of them, so this cost grows with $\K$, and it is paid early in the search.
Fixing a missing position needs a slot that is off and can be switched on at the right place.
When $\K$ is close to $\dist$, such spare slots run out as the search gets close to the target, and progress slows down.
With few evaluations, most of the effort goes into switching off wrong slots, so fewer slots are better.
With many evaluations, progress near the target is what limits the search, so spare slots pay off.

\subsection{Several Good Solutions}
\label{sec:theory:multi}

A multi-objective search looks for a whole front of good solutions, and they can lie at different distances from the rule.
We extend the model to $n_t$ targets at distances $\dist_1, \dots, \dist_{n_t}$, run NSGA-II on the distances to all targets,
and measure the mean, over targets, of the smallest distance reached.

The best $\K$ then follows the mean distance, not the smallest one.
With targets at distances 2 and 40 and 2000 evaluations, for example, the best $\K$ is 30, whereas for the near target alone it is 5.
Over all target sets we tried, the best $\K$ is up to 7.5 times the smallest distance with 465 evaluations and up to 15 times with 2000 evaluations.
A deviation rate chosen only to reach the nearest good solution is therefore too small for a front.

The best $\K$ also grows when unnecessary deviations are cheap.
In the base model, a deviation at a position where the target agrees with the rule costs one unit, like any other error.
If such a deviation cost only 0.1 instead, doubling $\K$ from 25 to 50 would no longer hurt at $\dist = 5$.

\subsection{Consequences for \RDD}
\label{sec:theory:consequences}

The analysis gives three guidelines, which the experiments test:
\begin{enumerate}
  \item With fewer slots than the rule distance, the best solutions are out of reach (Proposition~\ref{prop:floor}).
  Once the budget is large enough to get close to good solutions, $\K$ should exceed the typical rule distance.
  \item Larger budgets call for more slots (Table~\ref{tab:theory-budget}).
  \item In multi-objective search, the number of slots should be chosen for the mean distance to the front, not for the distance to the nearest good solution (Section~\ref{sec:theory:multi}).
\end{enumerate}
The rule distance of a real instance is unknown before the search.
This is why \RDD sets the deviation rate on training instances (Section~\ref{sec:method:k}).
